\begin{abstract}
  Advances in large language models (LLMs) have created new opportunities in data science, but their deployment is often limited by the challenge of finding relevant data in large data lakes. Existing methods struggle with this: both single- and multi-agent systems are quickly overwhelmed by large, heterogeneous files, and master–slave multi-agent systems rely on a rigid central controller that requires precise knowledge of each sub-agent’s capabilities, which is not possible in large-scale settings where the main agent lacks full observability over sub-agents’ knowledge and competencies. We propose a novel multi-agent paradigm inspired by the blackboard architecture for traditional AI models. In our framework, a central agent posts requests to a shared blackboard, and autonomous subordinate agents--either responsible for a partition of the data lake or retrieval from web--volunteer to respond based on their capabilities. This design improves scalability and flexibility by removing the need for a central coordinator to know each agent’s expertise or internal knowledge. We evaluate the approach on three benchmarks that require data discovery: KramaBench and modified versions of DSBench and DA-Code. Results show that the blackboard architecture substantially outperforms strong baselines, achieving 13\%–57\% relative improvements in end-to-end success and up to a 9\% relative gain in data discovery F1 over the best baseline. 
\end{abstract}